\chapter{Adaptadores de transporte}
\label{ch:transports}

\begin{note}
Stub. Tres transports delegan en \code{GraphHunterApi}: la app
desktop Tauri, el servidor HTTP basado en axum, y el servidor
MCP en TypeScript. Los tres están pensados para ser delgados
--- nada de lógica más allá de construir un request DTO y
renderizar la respuesta.
\end{note}

\section{Tauri}

Los comandos viven en
\filepath{app/src-tauri/src/commands/*.rs}. Cada cuerpo de
comando es $\leq$ 10 LOC:

\begin{lstlisting}[style=rust]
#[tauri::command]
pub fn cmd_run_hunt(
    api: State<Arc<GraphHunterApi>>,
    hypothesis: Hypothesis,
    time_window: Option<(i64, i64)>,
    max_results: Option<u32>,
) -> Result<RunHuntResponse, CommandError> {
    api.run_hunt(RunHuntRequest {
        session: None,
        hypothesis,
        time_window,
        max_results,
    })
    .map_err(CommandError::from)
}
\end{lstlisting}

Los comandos se registran en \filepath{app/src-tauri/src/lib.rs}
con la macro \code{tauri::Builder::invoke\_handler}.

\section{HTTP}

Los handlers viven en
\filepath{app/src-tauri/src/http\_api.rs} y comparten el mismo
\code{Arc<GraphHunterApi>} vía el trait \code{FromRef} de axum.
El puerto HTTP se elige al arrancar
(\code{util::resolve\_api\_port}) y se escribe al directorio de
config de la app Tauri para que el servidor MCP y la CLI puedan
encontrarlo.

Cada handler también es $\leq$ 10 LOC:

\begin{lstlisting}[style=rust]
async fn handler_run_hunt(
    State(api): State<Arc<GraphHunterApi>>,
    Json(body): Json<RunHuntRequest>,
) -> Response {
    match api.run_hunt(body) {
        Ok(v)  => ok_json(v),
        Err(e) => error_json(e),
    }
}
\end{lstlisting}

\section{MCP}

\filepath{graph-hunter-mcp/src/index.ts} expone $\sim$30
herramientas sobre el protocolo MCP. Cada handler de tool hace
una llamada HTTP al motor local vía
\code{http://127.0.0.1:<port>}; no enlaza el motor Rust en
absoluto. Este es el único transport que no sostiene un
\code{Arc<GraphHunterApi>} directo --- porque corre en otro
proceso (Node.js) --- pero ve exactamente la misma semántica
porque la capa HTTP es un wrapper delgado de la fachada.

\section{EventEmitter}

Los eventos de progreso (carga de sesión, ingesta, streaming
Sentinel) fluyen por un \code{Arc<dyn EventEmitter>} guardado en
\code{ApiState}. \code{TauriEventEmitter} puentea a
\code{AppHandle::emit}; \code{NoopEmitter} lo usan HTTP, MCP,
CLI y tests.

\begin{inthecode}
\filepath{app/src-tauri/src/commands/} --- módulos de comandos
Tauri.\\
\filepath{app/src-tauri/src/http\_api.rs} --- handlers HTTP.\\
\filepath{graph-hunter-mcp/src/index.ts} --- tools MCP.\\
\filepath{graph\_hunter\_api/src/events.rs} (o similar) ---
trait \code{EventEmitter}.
\end{inthecode}
